% !TeX program = xelatex
% !TeX encoding = UTF-8
% 2026 MathorCup 数学应用挑战赛——大数据竞赛
% 结构：封面 → 摘要 → 目录 → 正文 → 参考文献 → 附录
\documentclass{MathorCupBigData}

\usepackage{amsmath,amssymb,booktabs,graphicx,subcaption}
\usepackage{algorithm2e}
\usepackage{float}
\usepackage[section]{placeins}
\usepackage{longtable}
\usepackage{needspace}
\usepackage{tikz}
\usetikzlibrary{arrows.meta,positioning,shapes.geometric,fit,calc}

% ==== 竞赛信息（三处必填）====
\bianhao{MCB260XXXX}
\saidao{A}                 % A 或 B
\timu{在此填写论文题目}
\keyword{关键词1；关键词2；关键词3；关键词4}

\begin{document}

\makecover

% ==== 摘要（独立一页，建议 600–800 字，逐问一段，含方法名与数值结果）====
\begin{abstract}
针对问题一，……（说明数据预处理与建模方法，给出关键数值结果）。

针对问题二，……（说明模型与求解算法，给出指标与对比）。

针对问题三，……（说明优化或预测方案，给出结论量级）。

本文模型在验证集上的主要指标为……，相对基线提升……。
\end{abstract}

% ==== 目录（大数据赛官方模板保留目录）====
\tableofcontents
\clearpage

% ==== 正文 ====
\section{问题重述}
\subsection{问题背景}
\subsection{问题提出}

\section{问题分析}
\subsection{数据概况}
\subsection{总体思路}

\section{模型假设}
\begin{enumerate}
  \item 假设一：……
  \item 假设二：……
\end{enumerate}

\section{符号说明}
\begin{center}
\begin{tabular}{clc}
\toprule
符号 & 含义 & 单位 \\
\midrule
$X$ & 特征矩阵 & -- \\
$y$ & 目标变量 & -- \\
\bottomrule
\end{tabular}
\end{center}

\section{问题一模型的建立与求解}
\subsection{数据预处理}
\subsection{模型建立}
\subsection{求解与结果分析}

\section{问题二模型的建立与求解}
\subsection{模型建立}
\subsection{求解与结果分析}

\section{问题三模型的建立与求解}
\subsection{模型建立}
\subsection{求解与结果分析}

\section{模型评价与推广}
\subsection{模型优点}
\subsection{模型不足}
\subsection{模型推广}

% ==== 参考文献 ====
\begin{thebibliography}{99}
\bibitem{ref1} 作者. 题名[J]. 期刊, 年, 卷(期): 页码.
\bibitem{ref2} 作者. 书名[M]. 出版地: 出版社, 年.
\end{thebibliography}

% ==== 附录 ====
\appendix
\section{支撑材料清单}
\section{核心代码}
% \lstinputlisting[language=Python]{code/main.py}

\end{document}
